\section{Congruence quotients preserve relations, not ledgers}
\label{sec:quotients}

For a modulus $q\ge1$, let
\[
M_{r,q}=\{r^k\bmod q:k\ge0\},\qquad
X_{r,q}=(\Z/q\Z)\rtimes M_{r,q}.
\]
The labelled maps are
\begin{equation}
U_q(b,c)=(b+c,c),\qquad V_q(b,c)=(b,rc).
\label{eq:quotient-actions}
\end{equation}
When $V_q$ is not injective, reverse arcs remain formal reverses of edge
instances rather than inverse maps.

\begin{theorem}[Quotient boundary]
\label{thm:quotient}
For every $q\ge1$:
\begin{enumerate}
\item $U_qV_q=V_qU_q^r$, so the labelled affine relation word remains closed
and cyclically nonbacktracking in the formally symmetrized quotient;
\item the quotient also has the positive translation cycle
$U_q^q(0,1)=(0,1)$, primitive of length $q$ for $q\ge2$;
\item the affine relation polygon can lose vertex simplicity at small
moduli---for example at $(r,q)=(2,2)$;
\item if $q>r$, the relation polygon based at $(0,1)$ is vertex-simple, but
the quotient-created translation cycle still remains.
\end{enumerate}
Consequently the labelled relation descends, whereas the complete primitive
ledger does not.
\end{theorem}

\begin{proof}
Starting at $(b,c)$, both $V_q$ then $U_q$ and $U_q^r$ then $V_q$ end at
$(b+rc,rc)$.  Moreover $U_q^j(0,1)=(j\bmod q,1)$, so the first positive return
is $j=q$.  At $(r,q)=(2,2)$, the $V_q$ branch reaches multiplier zero, where
$U_q$ is a self-loop, while the $U_q^2$ branch already returns to its initial
vertex.  For $q>r$, residues $0,1,\ldots,r$ are distinct and the two branches
are internally disjoint.
\end{proof}

This result rules out a common shortcut.  A family of finite matrices may own
finite determinants, but it cannot be identified with the determinant of the
infinite positive Cayley graph without a separate convergence and descent
theorem.  The cycle $U_q^q$ is a concrete obstruction to ledger fidelity.
